\chapter*{Macromolecular Structures as Receptors}

The quality and accuracy of the results obtained from molecular docking, largely depends on the quality of the tertiary structure of the biological target. The structures are resolved using X-ray crystallography, Nuclear Magnetic Resonance (NMR) spectroscopy, and Cryo-Electron microscopy (CryoEM). These methods provide with high quality structural data, however, introduce errors such as missing atom information or multiple amino acid rotamers\cite{bib55, bib56}. Furthermore, X-ray crystallographic and CryoEM resolved structures are single snapshot of the receptor structure, while NMR derived structures provide insights in to the dynamics. Hence, it is crucial to evaluate and correct these structural defects before using them for molecular docking. Emphasis should be on the binding site amino acids, their protonation states, and appropriate rotameric states. Furthermore, the receptor's conformational state or functional states are important. Since most of the experimental structures do not contain this information, use of multiple structures resolved experimentally under different conditions or use MD simulations is highly recommended to include target flexibility. The worldwide protein Data Bank (wwPDB) is the primary source for receptor structures\cite{bib41}. In the absence of experimental data, computational methods, like template-based modelling and \textit{ab initio} predictions, can be used to predict the tertiary structures from primary sequences.  These theoretical models require thorough validation to ensure accuracy of the predicted tertiary folds. With recent developments in AI-based methods in protein structure prediction, such as AlphaFold\cite{bib37}, high quality structure are readily available in the wwPDB annotated as \href{https://pdb101.rcsb.org/learn/guide-to-understanding-pdb-data/computed-structure-models}{\textit{Computed Structure Models (CSMs)}}.

